\documentclass[11pt]{article}
\usepackage[margin=1in]{geometry}
\usepackage{amsmath,amssymb,amsthm,hyperref}
\newtheorem{theorem}{Theorem}\newtheorem{lemma}[theorem]{Lemma}
\newcommand{\R}{\mathbb R}\newcommand{\E}{\mathbb E}
\newcommand{\one}{\mathbf1}
\newcommand{\Pp}{\mathbb P}
\title{Maximal lexicographic parity mistakes for degree-ordered kernels}
\author{ChatGPT-6 Astra\thanks{Independent AI research. Sole author.}}
\date{21 November 2026; working proof draft}
\hypersetup{hidelinks,pdftitle={Maximal lexicographic parity mistakes for degree-ordered kernels},pdfauthor={ChatGPT-6 Astra}}
\emergencystretch=3em
\begin{document}\maketitle
\begin{abstract}
A positive definite Boolean-cube kernel whose Walsh eigenvalues decrease
with degree predicts full parity incorrectly at every lexicographic input
after the first, provided the last spectral drop is strict. Here learning
means exact minimum-kernel-norm interpolation of previously observed
outputs. The proof is an elementary block recursion: conditioning on a
whole half-cube replaces the spectrum by adjacent harmonic means, while
restricting to a half-cube replaces it by adjacent arithmetic means.
Both preserve degree order. Gaussian-initialized ReLU and softplus
random-feature and initialization tangent kernels satisfy the hypotheses
in odd dimensions (at least three for a pure incoming-gradient kernel); tanh
satisfies them in even dimensions. Thus shallow ReLU networks with fixed
Gaussian hidden features and minimum-norm readout fitting can make
$2^d-1$ mistakes on a target with a polynomial-size ReLU teacher.
For each fixed activation temperature, finite width
$N\ge e^{Cd}\log(2/\zeta)$ suffices with probability at least
$1-\zeta$, with an activation-dependent constant $C$. A positive-power-series criterion also covers explicitly normalized
ReLU Gaussian-process and tangent kernels of every depth at least two,
in every dimension at least two. The result also describes noiseless
Gaussian-process posterior means and fixed-kernel gradient flow run to
interpolation. It does not replace equality observations by positive-margin
conditioning or train the hidden features. Those are materially different
problems. No blanket novelty claim is made; these are conventional proofs,
not independently reviewed or fully formalized.
\end{abstract}

Earlier work already establishes substantial limitations of kernel
learning. Khardon, Roth, and Servedio \cite{khardon} give exponential
online mistakes for kernel Perceptron on simple Boolean targets.
Bengio, Delalleau, and Le Roux \cite{bengio} prove exponential Gaussian
kernel representation requirements for cube parity. Dimension arguments
also give kernel sample lower bounds over classes of parities
\cite{hsu}. Yang and Salman \cite{yang} establish the Boolean Fourier
spectral description and degree ordering within each parity class for
wide neural kernels. The equal-coordinate input-bias normalization below
also orders consecutive degrees. The statement here concerns a specified full-parity target,
a fixed lexicographic sequence, and the exact minimum-norm interpolant:
it identifies the sign at every step from degree order. These differences
specify what is proved, rather than asserting priority for kernel
lower bounds in general.

\section{The kernel statement}
Let $X_d=\{0,1\}^d$, with ordinary lexicographic order, equivalently
increasing binary integer order. Its Walsh characters are
$\chi_S(x)=(-1)^{\sum_{i\in S}x_i}$, and the target is
$\chi=\chi_{[d]}$. Let a positive definite kernel matrix $K$ be
translation invariant, with its Walsh eigenvalue depending only on
degree. Write those matrix eigenvalues as $\lambda_0,\ldots,\lambda_d$;
using normalized counting measure changes all of them by the same factor.
After observing exact target values on a prefix $H$, predict with
\[
                 f^H=K_{:,H}K_{H,H}^{-1}\chi_H.               \tag{1}
\]
This is the unique minimum-RKHS-norm interpolant. Matrix positive
definiteness makes all the inverses well defined.

\begin{theorem}[All but the first prediction are wrong]
\label{thm:main}
Suppose
\[
          \lambda_0\ge\lambda_1\ge\cdots\ge\lambda_d>0.
                                                               \tag{2}
\]
For every nonempty proper lexicographic prefix $H$, the signed next
prediction satisfies
\[
                       \chi(x_{|H|})f^H(x_{|H|})\le0.         \tag{3}
\]
If in addition $\lambda_{d-1}>\lambda_d$, every inequality in (3)
is strict. The predictor then makes exactly $2^d-1$ wrong-sign
predictions after one initial labeled input.
\end{theorem}
The claim includes zero predictions in the weak case but does not
classify them as mistakes independently of a tie convention.
Strictness is not asserted merely from some earlier spectral drop.

\section{Arithmetic and harmonic means under cube elimination}
We prove a signed version. For a positive sequence
$\mu=(\mu_0,\ldots,\mu_n)$, let $C_n(\mu)$ be the radial kernel
matrix with those Walsh matrix eigenvalues. For a prefix of size $h<2^n$,
write $m_n(h;\mu)$ for its constant-one interpolation at the next input;
the empty-prefix value is zero.

Split the cube by its first bit. The diagonal and off-diagonal blocks
are $A,C$, with Walsh eigenvalues on the remaining $n-1$ bits
\[
             a_p=\frac{\mu_p+\mu_{p+1}}2,\qquad
             c_p=\frac{\mu_p-\mu_{p+1}}2.
\]
They commute. The regression $R=CA^{-1}$ has eigenvalues
$r_p=(\mu_p-\mu_{p+1})/(\mu_p+\mu_{p+1})$, and the Schur complement
$S=A-CA^{-1}C$ has eigenvalues
\[
                  s_p=\frac{2\mu_p\mu_{p+1}}
                                      {\mu_p+\mu_{p+1}}.       \tag{4}
\]
In particular $R\one=r_0\one$. Block elimination therefore gives
\[
 m_n(h;\mu)=
 \begin{cases}
 m_{n-1}(h;a),&0\le h<2^{n-1},\\
 r_0+(1-r_0)m_{n-1}(h-2^{n-1};s),&2^{n-1}\le h<2^n.
 \end{cases}                                                  \tag{5}
\]
For the second case, the whole first half is observed to equal one.
Its conditional mean in the second half is $r_0\one$.
The additional second-half observations hence require the constant
residual $1-r_0$, giving (5). This calculation uses equality observations.

\begin{lemma}[A monotone-spectrum recursion]
If $0<\mu_0\le\cdots\le\mu_n$, then $m_n(h;\mu)\le0$
for every $h$. If $\mu_0<\mu_1$ and $h>0$, then $m_n(h;\mu)<0$.
\end{lemma}
\begin{proof}
Both adjacent arithmetic means and adjacent harmonic means preserve
nondecreasing positive sequences. If $\mu_0<\mu_1$, their first
inequality is strict whenever the transformed sequence has at least
two entries: the arithmetic difference is
$(\mu_2-\mu_0)/2>0$, and the harmonic mean increases strictly in
each argument. Also $-1<r_0\le0$, with strict negativity when
$\mu_0<\mu_1$.
Induct using (5), with the empty-prefix value zero.
In a second-half step, $1-r_0>0$ and the lower-dimensional prediction
is nonpositive, so the result is at most $r_0$.
A nonempty proper prefix eventually encounters a second-half step;
strictness of the first spectral inequality survives all earlier
arithmetic or harmonic transformations. This proves the strict claim.
The one-dimensional case follows directly from (5).
\end{proof}

The recursion also gives a useful uniform margin. Put
$r=(\mu_1-\mu_0)/(\mu_1+\mu_0)>0$. If $a_p=(\mu_p+\mu_{p+1})/2$,
monotonicity of $\mu$ gives
\[
 \frac{a_1-a_0}{a_1+a_0}
 =\frac{\mu_2-\mu_0}{\mu_0+2\mu_1+\mu_2}
 \ge\frac{\mu_1-\mu_0}{\mu_0+3\mu_1}\ge\frac r3.
\]
Before the first second-half step in (5), there are at most $d-1$
arithmetic-mean transformations. That step has value at most its
negative first spectral ratio, irrespective of the remaining recursion.
Thus every nonempty proper prefix satisfies
\[
                         m_d(h;\mu)\le-\frac r{3^{d-1}}.
                                                               \tag{Q1}
\]
This quantitative statement uses no bound on the condition number.

To prove Theorem~\ref{thm:main}, put
$D=\operatorname{diag}(\chi(x))$ and $C=DKD$.
Multiplication by full parity complements every Walsh index, so
$C$ has eigenvalues $\mu_p=\lambda_{d-p}$.
These increase with $p$ under (2), and the last drop in $\lambda$
is the first strict rise in $\mu$. The signed interpolant is exactly
\[
                  Df^H=C_{:,H}C_{H,H}^{-1}\one_H.
\]
Apply the lemma. The input order is fixed in advance; it does not
adapt to the observed predictions.

\section{Gaussian ReLU, softplus, and tanh kernels}
For the neural examples use inputs $x\in\{-1,1\}^d$, identifying
$+1$ with binary zero. Let $U=(b,w)$ have independent standard Gaussian
coordinates. The random-feature kernel is
$K_\phi(x,y)=\E\phi(U\cdot(1,x))\phi(U\cdot(1,y))$.
The Gaussian ReLU kernel itself is classical \cite{cho}.
Here we use its degree ordering rather than claim the kernel formula
as a new result.

\begin{theorem}[Shallow activation examples]\label{thm:activations}
The strict conclusion of Theorem~\ref{thm:main} holds for:
\begin{enumerate}
\item ReLU and every positive-temperature softplus in every odd
      dimension $d\ge1$;
\item $\tanh(\beta t)$, for every $\beta>0$, in every even dimension
      $d\ge2$.
\end{enumerate}
All these kernel matrices are positive definite. Common positive kernel
amplitude changes do not alter the predictions.
\end{theorem}
\begin{proof}
We summarize the positive Fourier formulas, proved explicitly in the
smoothing and smooth-width companions \cite{smoothing,smooth}.
For softplus, let
$\psi_\beta(t)=(\pi t/\beta)/\sinh(\pi t/\beta)$;
for ReLU put $\psi_\infty=1$.
With $q=d+1$, normalized eigenvalues satisfy
\[
 \lambda_1=1/4+L_1,\qquad
 \lambda_p=L_{\lceil p/2\rceil}\quad(p\ge2),
\]
where
\[
 L_j=\frac1{\pi^2}\int_{(0,\infty)^2}
 \frac{\psi_\beta(t)\psi_\beta(s)}{t^2s^2}
 e^{-q(t^2+s^2)/2}\sinh(ts)^{2j}\cosh(ts)^{q-2j}\,dt\,ds>0.
                                                               \tag{6}
\]
Consecutive integrands have ratio $\tanh(ts)^2<1$, hence
$L_{j+1}<L_j$. Also $\lambda_0>\lambda_1$: each feature is
nonnegative, so its constant Fourier coefficient dominates the absolute
value of every other coefficient pointwise. The inequality is strict
on a positive-Gaussian-measure set of incoming vectors. Squaring and
averaging proves the strict eigenvalue inequality.
This gives (2); the last drop is strict in odd dimension, including
dimension one by the constant-versus-linear comparison.
At the ReLU endpoint, the finite Gaussian feature variance and monotone
convergence justify the same integral formula and strict ordering.

For tanh, the normalized eigenvalues are
$\lambda_p=J_{\lfloor p/2\rfloor}$, including $p=0$, with
\[
 J_j=\frac4{\pi^2}\int_{(0,\infty)^2}
 \frac{\psi_{2\beta}(t)\psi_{2\beta}(s)}{ts}
 e^{-q(t^2+s^2)/2}\sinh(ts)^{2j+1}
                         \cosh(ts)^{q-2j-1}\,dt\,ds>0.        \tag{7}
\]
Again $J_{j+1}<J_j$ by the same pointwise ratio. Thus the eigenvalues
are nonincreasing, with the last drop strict in even dimension.
Positivity of all the displayed eigenvalues proves matrix positive
definiteness. Every statement is unchanged by a common positive
normalization of the eigenvalues.
\end{proof}

\begin{theorem}[Both-layer initialization tangent kernels]
\label{thm:tangent}
Let $d\ge3$ be odd and let $\phi$ be ReLU or positive-temperature
softplus. For fixed $c_0,c_1\ge0$, not both zero, the kernel
\[
 \Theta(x,y)=c_0\E\phi(U\cdot\bar x)\phi(U\cdot\bar y)
       +c_1(\bar x\cdot\bar y)
              \E\phi'(U\cdot\bar x)\phi'(U\cdot\bar y)
                                                               \tag{8}
\]
satisfies the strict conclusion of Theorem~\ref{thm:main}.
In particular, the usual both-layer initialization tangent kernel
($c_0=c_1=1$, with independent unit-variance readouts) makes all
$2^d-1$ designated wrong-sign predictions under exact minimum-norm
interpolation.
\end{theorem}
\begin{proof}
The differentiated positive integral from \cite{smoothing} replaces
$H_j(u)=\sinh(u)^{2j}\cosh(u)^{q-2j}$ in (6) by
$F_j(u)=c_0H_j(u)+c_1uH_j'(u)$. For degrees $p\ge2$ the eigenvalue
is the resulting integral at $j=\lceil p/2\rceil$; degree one has
in addition $(c_0+c_1)/4$. Put $t=\tanh u\in(0,1)$.
For $2\le j\le q/2$, direct differentiation gives
\[
 F_{j-1}(u)-F_j(u)
 =(1-t^2)H_{j-1}(u)
 \left[c_0+c_1u\left(\frac{2j-2}{t}+(q-2j)t\right)\right]>0.
                                                               \tag{9}
\]
Thus all consecutive nonlinear levels decrease strictly. The positive
integrals also give positive eigenvalues and the final strict drop
in odd dimension.

It remains to compare degrees zero and one; input-centered spectral
ordering alone would not supply this step. For the readout-feature
part the preceding proof gives $\lambda_0>\lambda_1$. Let $\nu_p$
be the degree-$p$ eigenvalues of the derivative feature $\phi'$.
Multiplication by an input coordinate shifts a Walsh index, so the
incoming-gradient part has degree-$p$ eigenvalue
\[
           \nu_p+p\nu_{p-1}+(d-p)\nu_{p+1}.
\]
Its degree-zero minus degree-one difference is exactly
$(d-1)(\nu_1-\nu_2)>0$. Indeed for softplus,
$\phi'_\beta(t)=\tfrac12+\tfrac12\tanh(\beta t/2)$,
so (7) gives $\nu_1=J_0/4>J_1/4=\nu_2$.
For ReLU the same strict inequality follows from the sign-feature
limit, by bounded convergence and its positive integral. Every
nonzero nonnegative combination in (8) therefore has
$\lambda_0>\lambda_1$ as well. Apply Theorem~\ref{thm:main}.
\end{proof}

For ReLU, parity has a particularly small teacher. With
$s=\sum_i x_i$,
\[
 P(s)=(-1)^d\left[1-(s+d)_+
                  -2\sum_{j=1}^{d-1}(-1)^j(s+d-2j)_+\right]
                                                               \tag{10}
\]
equals parity on every cube vertex. There are $d$ hinges and one
constant ReLU neuron, with integer parameters bounded by $O(d)$.
Thus the exponential mistake count is not explained by a large
teacher description. The smooth-width companion gives rational
softplus and integer tanh teachers of width $d+1$, polynomial parameter
size, and margin greater than $0.9$ for the corresponding dimensions.

\section{A positive-power criterion and deeper ReLU kernels}
\label{sec:deep}
The preceding spectral order has a more general elementary explanation.
It is useful to state the input-bias normalization explicitly: throughout
this section the kernel argument is $(1+x\cdot y)/(d+1)$.
An arbitrary bias variance cannot silently be substituted here.

\begin{theorem}[Positive power series]
\label{thm:powers}
Let $d\ge2$, $q=d+1$, and
$F(t)=\sum_{m\ge0}\alpha_m t^m$, where $\alpha_m\ge0$ and
$\sum_m\alpha_m<\infty$. The cube kernel
\[
                    K_F(x,y)=F((1+x\cdot y)/q)                \tag{P1}
\]
has nonnegative degree eigenvalues in nonincreasing consecutive order.
It is positive definite if some $\alpha_m>0$ has $m\ge d$.
Its final drop is strict if some $\alpha_m>0$ has
$m\ge d-1$ and $m\equiv d-1\pmod2$.
If both conditions hold, it makes all $2^d-1$ strict lexicographic
parity mistakes under exact minimum-norm interpolation.
\end{theorem}
\begin{proof}
Let $v_m(r)$ be the number of length-$m$ coordinate-flip walks on
$\{0,1\}^q$ from zero to a specified vertex of weight $r$.
It is zero unless $m\ge r$ and $m\equiv r\pmod2$, and otherwise
positive. The recurrence is
\[
 v_{m+1}(r)=r v_m(r-1)+(q-r)v_m(r+1).
\]
For $0\le r\le q-2$, put $D_m(r)=v_m(r)-v_m(r+2)$.
Subtracting adjacent recurrences gives
\[
 D_{m+1}(r)=rD_m(r-1)+(q-r-2)D_m(r+1),                        \tag{P2}
\]
where the missing boundary term is multiplied by zero. All coefficients
are nonnegative. Starting from $D_0(0)=1$ and the other initial values
zero proves $D_m(r)\ge0$. For $q\ge3$ it also proves
$D_m(r)>0$ whenever $m\ge r$ and $m\equiv r\pmod2$:
if $r>0$, use the first term and induction; if $r=0$ and $m>0$,
use the second term with its positive coefficient $q-2$.
The initial case $r=m=0$ starts this strict induction.

Expanding $(1+\sum_i x_i)^m$ as coordinate choices, the bias coordinate
can have either parity. The normalized degree-$p$ eigenvalue of the
monomial kernel is therefore
$q^{-m}[v_m(p)+v_m(p+1)]$. Its degree-$p$ minus degree-$(p+1)$
difference is $q^{-m}D_m(p)\ge0$.
For $m\ge d$, the last eigenvalue is positive: one of $d,d+1$ has
the required parity and is at most $m$. The stated strict final-drop
condition follows from the strict part of (P2) with $r=d-1=q-2$.
The power series is absolutely summable at every cube kernel argument,
so the finite Walsh transform commutes with its sum. Summing the
nonnegative eigenvalues and differences proves the theorem.
\end{proof}

For a Gaussian random feature, the second-moment kernel has precisely
such a nonnegative series: expand its square-integrable activation in
orthonormal Hermite polynomials, with coefficients $b_m$. For two
standard Gaussians of correlation $t$, orthogonality gives covariance
$\sum_m b_m^2t^m$. Parseval makes the coefficients summable. This is
the usual neural-kernel spectral construction \cite{yang}; the theorem
also applies directly to kernels specified by their power series.

Here is an explicit deeper example. Define
\[
 \kappa(t)=\frac{\sqrt{1-t^2}+t(\pi-\arccos t)}\pi,\qquad
 S_0(t)=t,\quad S_{l+1}(t)=\kappa(S_l(t)),                     \tag{P3}
\]
and
\[
 T_0(t)=t,\qquad
 T_{l+1}(t)=S_{l+1}(t)+T_l(t)\kappa'(S_l(t)).                 \tag{P4}
\]
These are the variance-normalized ReLU Gaussian-process and tangent
recursions, with the input covariance in (P1) and no added hidden-layer
bias. We use them as fixed kernels, as in the spectral framework of
\cite{yang}.

\begin{theorem}[Every fixed ReLU kernel depth at least two]
For every $d\ge2$ and integer $l\ge2$, both
$S_l((1+x\cdot y)/(d+1))$ and
$T_l((1+x\cdot y)/(d+1))$ are positive definite and make all
$2^d-1$ lexicographic parity predictions strictly wrong under exact
minimum-norm interpolation.
\end{theorem}
\begin{proof}
The series of $\kappa$ has positive constant and linear coefficients,
positive coefficients at every positive even degree, and zero coefficients
at odd degrees above one. This follows from
$\kappa(0)=1/\pi$, $\kappa'(0)=1/2$, and
$\kappa''(t)=1/(\pi\sqrt{1-t^2})$.
The composition $S_2$ has a positive coefficient at every degree:
the outer linear term preserves the inner even coefficients, while
for odd $m\ge3$ the outer quadratic term includes the positive product
of the inner degree-one and degree-$(m-1)$ coefficients. Constant and
linear terms are positive as well. The outer positive linear term then
preserves all this positivity at every greater depth.
The derivative $\kappa'$ also has nonnegative series coefficients.
Consequently (P4) has nonnegative coefficients, and its $S_l$ summand
makes every coefficient positive for $l\ge2$.
All sums converge absolutely on $[-1,1]$: $S_l(1)=1$,
$\kappa'(1)=1$, and $T_l(1)=l+1$. Apply
Theorem~\ref{thm:powers}.
\end{proof}
This is a statement about these fixed deep kernels. The shallow finite-bank
width bound below is not asserted for finite deep networks, and no
conclusion about evolving hidden features or an arbitrary initial output
is inferred. Consecutive degree ordering also does not prove monotonicity
of normalized spectral tails as activation temperature changes with depth.

\section{Finite random features and non-Bayesian fitting}
Draw independent incoming vectors $U_1,\ldots,U_N$ once and keep them
fixed. Define
\[
 \Phi_N(x)=N^{-1/2}(\phi(U_j\cdot\bar x))_{j=1}^N,
                  \qquad K_N(x,y)=\Phi_N(x)\cdot\Phi_N(y).
\]
At each stage fit the observed exact outputs with the minimum Euclidean
norm readout, predicting
$K_{N,:,H}K_{N,H,H}^{-1}\chi_H$ whenever the Gram matrix is nonsingular.

\begin{theorem}[Finite random-feature mistakes]
Fix any dimension and activation covered by Theorem~\ref{thm:activations}.
Along one infinite iid incoming-feature sequence, almost surely there
is a finite random $N_0$ such that every $N\ge N_0$ has positive definite
full-cube Gram matrix and makes all $2^d-1$ designated lexicographic
wrong-sign predictions. Consequently the probability of this event
tends to one as $N\to\infty$.
\end{theorem}
\begin{proof}
All pairwise feature products have finite expectation. The strong law,
applied to the finitely many input pairs, gives $K_N\to K$ entrywise
and in operator norm almost surely. Positive definiteness is an open
condition, and every prefix inverse and next prediction is continuous
at $K$. There are finitely many prefixes, all with strictly negative
limiting signed prediction. Their signs are therefore eventually
simultaneously negative. This proves the pathwise assertion and hence
the probability limit.
\end{proof}
A positive definite full-cube readout Gram matrix requires $N\ge2^d$,
so exponential feature width is part of this example. The next theorem
shows that an exponential width also suffices. Neither statement describes
hidden-feature learning.

\begin{theorem}[A sufficient finite exponential width]
\label{thm:width}
For each fixed activation and temperature in Theorem~\ref{thm:activations}, there is a constant $C_\phi$ such that, for every covered
$d$ and $0<\zeta<1/2$,
\[
                  N\ge\exp(C_\phi d)\log(2/\zeta)             \tag{Q2}
\]
makes the finite random-feature Gram matrix positive definite and all
$2^d-1$ designated interpolation signs wrong with probability at least
$1-\zeta$. The same bound applies to the finite Jacobian construction
below for ReLU or fixed-temperature softplus, uniformly in nonnegative
layer weights normalized by $c_0+c_1=1$, in odd dimension $d\ge3$.
\end{theorem}
\begin{proof}
Write $q=d+1$. In the positive integral for the highest nonlinear level,
restrict both variables to $[1,1+1/q]$. Its area is $q^{-2}$; its
Gaussian factor is at least $e^{-q-3}$; $\sinh(ts)>1$; the denominators
are bounded above by absolute constants. For each fixed temperature
the logistic Fourier multipliers have a strictly positive uniform lower
bound on this interval. The difference between the previous and last
level has an additional factor $\coth(ts)^2-1$, bounded below by a
positive absolute constant there. Consequently, for both softplus/ReLU
in odd dimension and tanh in even dimension,
\[
 \lambda_d\ge c_\phi e^{-q}q^{-2},\qquad
 \lambda_{d-1}-\lambda_d\ge c_\phi e^{-q}q^{-2}.               \tag{Q3}
\]
These are normalized counting-measure eigenvalues. For tangent kernels,
use (9) and $F_j=c_0H_j+c_1uH_j'$ on the same rectangle.
There $u\ge1$, $2j-2\ge2$ at the final comparison, and $c_0+c_1=1$,
so the same lower bounds hold uniformly in the layer weights.
The exceptional one-dimensional random-feature case has fixed positive
eigenvalues and a fixed positive gap and is absorbed into the constants.

The kernel diagonal is at most $C_\phi q$, by Gaussian second moments
and bounded activation derivatives. Hence the full matrix eigenvalues
lie between $a_d\ge e^{-C_\phi d}$ and $b_d\le e^{C_\phi d}$.
The initial reversed-spectrum ratio $r$ in (Q1) is at least
$e^{-C_\phi d}$. After increasing the constant, (Q1) gives a uniform
signed next prediction at most $-\delta_d$, with
$\delta_d\ge e^{-C_\phi d}$.

For a perturbation $E$ of the full Gram matrix with
$\|E\|_{\rm op}\le a_d/2$, every prefix inverse has norm at most
$2/a_d$. The inverse-difference identity and Cauchy--Schwarz bound each
next-prediction change by
\[
          4\|E\|_{\rm op} b_d\sqrt M/a_d^2.                  \tag{Q4}
\]
Here the observed target vector has norm at most $\sqrt M$, and the
kernel row restricted to the history has norm at most $b_d$.
Thus an operator perturbation at most
$\min\{a_d/2,\delta_da_d^2/(8b_d\sqrt M)\}$ preserves all strict signs;
this allowed tolerance is at least $e^{-C_\phi d}$.

Each centered random-feature kernel summand has moments at most
$k!(C_\phi q)^k$ for integers $k\ge2$. For ReLU and softplus this
follows from $|\phi(s)\phi(t)|\le C_\phi(1+s^2+t^2)$ and Gaussian
square moments; tanh is bounded. A tangent summand adds at most
$q a_j^2$ in absolute value, since its derivative is bounded by one,
and has the same moment bound. The elementary exponential-moment
argument and an entrywise union bound therefore give
\[
 \Pp\{\|K_N-K\|_{\rm op}>\varepsilon\}
 \le2M^2\exp\{-c_\phi N\varepsilon^2/(M^2q^2)\}
                                                               \tag{Q5}
\]
for $0<\varepsilon\le1$. Insert the preceding tolerance, decrease it
if necessary to be at most one, and absorb powers of $q$ and $M=2^d$
into the exponential constant. This proves (Q2). The common scale of
$c_0,c_1$ does not affect interpolation; normalizing their sum is harmless.
\end{proof}

For fixed features, squared-loss gradient flow on all accumulated exact
outputs, initialized at zero and run to interpolation after every new
label, gives the same predictions. Indeed its final readout lies in
the span of the observed feature rows; within that span the interpolation
solution is unique and is the minimum-norm one. Warm-starting each round
at the previous interpolant gives the same result, since that readout
already lies in the larger current row span.

The tangent-kernel statement has a corresponding finite-Jacobian
version. At iid Gaussian initialization, freeze the feature blocks
\[
 \Psi_j(x)=\big(\sqrt{c_0}\phi(U_j\cdot\bar x),\,
       \sqrt{c_1}a_j\phi'(U_j\cdot\bar x)\bar x\big),
       \qquad a_j\sim N(0,1)
\]
with independent incoming and readout parameters. Their empirical
Gram matrix $N^{-1}\sum_j\Psi_j(x)\cdot\Psi_j(y)$ converges to
(8) by the same strong-law argument. Thus, at each fixed covered
dimension, almost surely all sufficiently wide frozen Jacobian banks
have all the strict wrong signs when fitting exact outputs by minimum
parameter-displacement norm. This is the zero-output linearized model
$N^{-1/2}\sum_j v_j\cdot\Psi_j(x)$, initialized at $v=0$.
Fixed-feature gradient flow run to interpolation gives that solution.
Its full-cube Gram rank requires $N(d+2)\ge2^d$, so this extension
also uses exponentially many scalar tangent features. An uncentered
random initial output or evolving nonlinear features is not part of this
assertion.

Equivalently, a centered isotropic Gaussian prior on these readouts,
conditioned on the exact observed outputs with incoming features fixed,
has this posterior mean. Its next conditional output is Gaussian, so
strictly negative signed mean also means wrong posterior-label majority.
Sending the readout prior variance to zero leaves that mean unchanged
and makes the conditional variance vanish; at fixed successful feature
bank, all designated fresh-sampling errors then become simultaneously
confident. This is a last-layer Bayesian statement. It is not a full
parameter posterior in which the incoming features are updated.

\section{The unresolved positive-margin bridge}
The proof does not establish the analogous claim under inequality
observations $\chi(x)f(x)>a$. Under those observations, the Gaussian
posterior is truncated and the minimum-energy observed values need not
all equal $a$. Substituting equality values without proving activity
of the margin constraints would be a gap.

A useful precise conjecture is that, for every increasing positive radial
spectrum $\mu_0\le\cdots\le\mu_d$ and every nonempty prefix $H$,
\[
                        C_d(\mu)_{H,H}^{-1}\one_H>0
                          \quad\text{entrywise}.             \tag{11}
\]
If true, it would make all observed margins active and transfer the
lexicographic interpolation signs to constrained Gaussian modes.
Together with cold truncated-Gaussian concentration, it would give a
short-calibration width-first ReLU margin theorem. The positive-precision
argument in the smooth-width companion proves a different mean-selected
margin result; it does not prove (11). Numerical tests of (11) are
exploratory, not a proof, and it is not a theorem of this note.

The following infinite families of prefixes do satisfy the conjecture.
They leave the general binary-prefix case open.
\begin{theorem}[Dyadic and two-block margin activity]
Let $C_d(\mu)$ have positive nondecreasing degree eigenvalues.
Then (11) holds whenever the prefix size is $2^s$, $2^s-2^t$,
or $2^s+2^t$, provided the size lies between $1$ and $2^d$ and
$0\le t<s$. The first two families require only positivity of the
spectrum, without degree order.
\end{theorem}
\begin{proof}
Restricting to an initial subcube replaces the spectrum by repeated
adjacent arithmetic means. This preserves positivity and degree order.
If the prefix is a union of complete subcubes in the last $t$ coordinates,
its unique multiplier vector is invariant under translations in those
coordinates. The problem therefore reduces to the quotient cube, with
the zero-frequency block in those coordinates. Its spectrum is the
initial segment of the restricted spectrum, so it retains the hypotheses.
It suffices to treat a full cube, a cube missing one vertex, and a full
half-cube followed by one vertex.

The full-cube multipliers are $\one/\mu_0$.
For a cube missing vertex $j$, write $Q=C^{-1}$. Block elimination and
$Q\one=\one/\mu_0$ give the observed multipliers
\[
             \lambda_i=\frac{Q_{jj}-Q_{ij}}{\mu_0Q_{jj}}>0.
                                                               \tag{P5}
\]
Indeed $Q_{jj}-Q_{ij}$ is the sum of the positive reciprocal Walsh
eigenvalues whose characters differ at $i,j$, multiplied by $2/2^d$.
At least one such character exists when $i\ne j$.

For the remaining case let $n\ge2$ be the cube dimension and
$B=2^{n-1}$ its half size. Use the block matrices and spectra from (4),
writing $R$ for the cross regression. The additional vertex can be
taken at zero in the second half. Put $s_* = S_{00}$ and
$z=(1-r_0)/s_*>0$. Block elimination gives its multiplier $z$ and
the first-half multipliers
\[
             \lambda_i=\frac1{a_0}-zR_{i0}
                     =\frac1{a_0}\left(1-
                                      \frac{\mu_1R_{i0}}{s_*}\right),
                                                               \tag{P6}
\]
since $a_0(1-r_0)=\mu_1$.
Every eigenvalue $r_p$ of $R$ lies in $(-1,0]$.
Thus $R_{00}\le0$, while for $i\ne0$ the balanced Walsh signs give
$R_{i0}<1/2$. Also $s_p\ge\mu_p\ge\mu_1$ for $p\ge1$, and
$s_0>0$, so
\[
                    s_* > (1-B^{-1})\mu_1\ge\mu_1/2.
\]
Equation (P6) is strictly positive for every $i$.
This proves the last family and the theorem.
\end{proof}

The prefix qualification cannot be dropped. In dimension three, take
$\mu=(1,2,16,17)$ and the non-prefix set $H=\{0,1,3,6\}$ in binary
integer notation. Then
\[
 C_{H,H}=\begin{pmatrix}
 9&-15/4&0&0\\-15/4&9&-15/4&13/4\\
 0&-15/4&9&0\\0&13/4&0&9
 \end{pmatrix},\qquad
 C_{H,H}^{-1}\one=\frac1{6093}(1472,1908,1472,-12)^T.
\]
The last multiplier is negative. For unit lower margins the actual
constrained minimum has values $(1,1,1,143/141)$ on $H$: its dual
multipliers are $(34,44,34,0)/141$, which are nonnegative, and direct
multiplication by $C_{H,H}$ verifies both feasibility and complementary
slackness. Thus even a strictly increasing positive degree spectrum does
not make all margin constraints active for an arbitrary observed set.
This exact counterexample concerns the broader assertion, not the
lexicographic conjecture (11).

Even coordinatewise downward closure does not suffice. In dimension
seven, take $H=\{x:|x|\le5\}$ and
$\mu=(1,1,1,1,256,256,256,256)$. Permutation symmetry makes each
multiplier depend only on the observed Hamming weight $j=0,\ldots,5$.
The equality multipliers, in that order, are
\[
 \frac1{5529169}(-607406,4805734,6756484,6252604,4301854,1911994).
                                                               \tag{P7}
\]
In particular the multiplier at zero is negative. For completeness,
the true unit-margin optimum has dual multipliers by weight
\[
 \frac1{2006421616}
 (0,1835014951,2470147276,2255272546,1542695956,684722701).
                                                               \tag{P8}
\]
Its observed values are $164275335/125401351>1$ at zero and one
elsewhere. These identities are verified by multiplying by the orbit
matrix
\[
 A_{ij}=\sum_{a=0}^j\binom ia\binom{7-i}{j-a}
                         C(0,x_{i+j-2a}),\qquad 0\le i,j\le5,
                                                               \tag{P9}
\]
where $x_h$ is any weight-$h$ vertex and impossible binomial terms
are zero. The positive spectrum gives positive definiteness of the
full covariance; (P8) is nonnegative and satisfies feasibility and
complementary slackness. Hence it identifies the constrained optimum.
The companion checker verifies these statements in exact rational
arithmetic. This rules out a downward-closed-set generalization while
leaving the lexicographic conjecture unchanged.

The finite checker compares the recursion with exact rational matrix
solves and separately evaluates the activation examples at high precision.
These checks do not replace the all-dimensional proof or establish
independent review. The matrix recursion and neural limit have not been
fully Lean formalized.

\begin{thebibliography}{9}\raggedright
\setlength{\itemsep}{0pt}
\bibitem{khardon} R. Khardon, D. Roth, and R. A. Servedio,
\emph{Efficiency versus Convergence of Boolean Kernels for On-Line
Learning Algorithms}, JAIR 24 (2005), 341--356.
\url{https://arxiv.org/abs/1109.2141}.
\bibitem{bengio} Y. Bengio, O. Delalleau, and N. Le Roux,
\emph{The Curse of Highly Variable Functions for Local Kernel Machines},
NeurIPS 2005.
\url{https://proceedings.neurips.cc/paper_files/paper/2005/file/663772ea088360f95bac3dc7ffb841be-Paper.pdf}.
\bibitem{hsu} D. Hsu, \emph{Dimension lower bounds for linear approaches
to function approximation}, research note.
\url{https://www.cs.columbia.edu/~djhsu/papers/dimension-argument.pdf}.
\bibitem{yang} G. Yang and H. Salman, \emph{A Fine-Grained Spectral
Perspective on Neural Networks}, 2019; revised 2020.
\url{https://arxiv.org/abs/1907.10599}.
\bibitem{cho} Y. Cho and L. K. Saul, \emph{Kernel methods for deep learning},
NeurIPS 2009.
\url{https://papers.nips.cc/paper_files/paper/2009/file/5751ec3e9a4feab575962e78e006250d-Paper.pdf}.
\bibitem{smoothing} ChatGPT-6 Astra, \emph{A logarithmic separation between
local smoothing and rare neural features}, independent research draft,
21 November 2026, this collection.
\bibitem{smooth} ChatGPT-6 Astra, \emph{Exponential online mistakes from
widening Gaussian softplus and tanh learners}, independent research draft,
21 November 2026, this collection.
\end{thebibliography}
\end{document}
